\section{Delimitaciones y limitaciones del trabajo}
\label{sec:delimitaciones-limitaciones}

% ~150-200 palabras, entre 3 y 5 bullets en total (según plantilla
% institucional). La plantilla distingue dos ideas distintas:
%   - Delimitaciones: aspectos que NO serán cubiertos durante ni al final
%     de la investigación (decisión propia de alcance).
%   - Limitaciones: aspectos que pudieran afectar la investigación pero
%     que están fuera del control del investigador.
% También hay que analizar el entorno donde se pretende validar el
% producto de la investigación.

\subsection{Delimitaciones}
\begin{itemize}
  \item La orquestación se implementará mediante herramientas-agente de codificación que accederán a modelos alojados en la nube a través de un enrutador de modelos (por ejemplo, OpenRouter), descritas en la Sección~\ref{sec:metodologia}, lo que permitirá emplear modelos de distintos proveedores sin requerir infraestructura local; en consecuencia, la evaluación se limitará a modelos accesibles en la nube y no contemplará despliegues en infraestructura propia.
  \item La evaluación se limitará a la corrección funcional y la estructura arquitectónica del código generado, por lo que no se considerarán otros aspectos de calidad del software como el rendimiento o la seguridad.
  \item La comparación se realizará en un entorno controlado de desarrollo sobre un lenguaje de programación específico (TypeScript), no en un entorno de producción.
\end{itemize}

\subsection{Limitaciones}
Las principales limitaciones de este trabajo se analizan como amenazas a la validez:
\begin{itemize}
  \item \textbf{Validez interna.} La condición multiagente difiere de la línea base en tres dimensiones simultáneas (número de agentes, heterogeneidad de modelos y especialización de roles), por lo que no es posible aislar cuál de ellas explica las diferencias observadas en la corrección funcional y la estructura del código. Para controlar la calidad del modelo, la línea base emplea el modelo individual de mayor capacidad (el mismo que ocupa el rol de arquitecto), de modo que el modelo más fuerte esté presente en ambas condiciones y la comparación se haga contra una línea base deliberadamente exigente. Aun así, dentro del sistema multiagente cada rol usa un modelo distinto, por lo que no es posible distinguir si un resultado se debe al rol o al modelo que le tocó (el arquitecto, por ejemplo, siempre lleva el modelo más capaz). Este estudio asume esa limitación y no separa ambos efectos. Además, los modelos de lenguaje son no deterministas y su desempeño puede variar entre ejecuciones e incluso en el tiempo por actualizaciones de los proveedores; para acotar este efecto se realizarán K repeticiones por tarea y condición y se fijarán los modelos y sus versiones durante toda la experimentación.
  \item \textbf{Validez externa.} La evaluación se limitará a un único dominio (el backend de un sistema de gestión de negocio en TypeScript), a un conjunto acotado de tareas y a modelos accesibles en la nube, por lo que sus resultados no son directamente generalizables a otros dominios, lenguajes o entornos de desarrollo. Asimismo, el costo de las ejecuciones limita el número de tareas, repeticiones y modelos que pueden evaluarse, y la fase de piloto emplea modelos de capacidad media, lo que restringe aún más el alcance de las conclusiones.
  \item \textbf{Validez de constructo.} Las métricas empleadas son aproximaciones de los conceptos que pretenden medir: el \textit{pass rate} (la proporción de pruebas automatizadas que aprueba el código) depende de la calidad de las suites de prueba y, al calcularse solo sobre el código que compila, se interpretará junto con la compilabilidad; la complejidad cognitiva, por su parte, es una medida indirecta de la estructura del código. Además, el diseño de los \textit{prompts} influye en el código generado; como la estructura de los \textit{prompts} difiere necesariamente entre el sistema multiagente y el agente único, lo que se mantendrá idéntico entre condiciones será la especificación de cada tarea (casos de uso y criterios de aceptación) y el esfuerzo de ajuste invertido en los \textit{prompts}, cuyo diseño definitivo se congelará tras el piloto.
\end{itemize}
